\chapter{Threats to validity}
\label{ch:threats}

This chapter focuses on the limitations that matter most for interpreting the results in Chapter~\ref{ch:results}. Some mitigations have already been described where the relevant design choices are introduced, such as the three-way inter-rater check, deterministic ranking tie-breaks, the multi-knob sensitivity analysis, and reproducible notebook sources for the reported numbers. They are not repeated here unless they affect the remaining threat directly.

\section{Construct and weight calibration}
\label{sec:threats-construct}

\paragraph{There is no external ground truth for ``refactoring quality''.}
The process-quality score $J$ measures behaviours that the literature treats as signs of a weaker refactoring process, such as broken builds, skipped tests, manual edits where an IDE refactoring would have worked, and long gaps between commits. It should not be read as a predictor of bug rates, time-to-merge, or other downstream outcomes, and Chapter~\ref{ch:methodology} does not make that claim. Because there is no dataset linking session-level process signals to downstream maintenance outcomes at this granularity, the weights are grounded in the literature-supported severity orderings in \S\ref{sec:methodology-weights}, rather than fitted to outcome data.

\paragraph{The weights are plausible, but not uniquely calibrated.}
The production weights follow the severity ordering supported by the literature, but the exact values are still a design choice. Other values within the same ordering could also be justified. In \S\ref{sec:results-sensitivity}, top-$1$ stability is $96.5\%$ under single-knob perturbation and $84.6\%$ under multi-knob Monte Carlo on the rankable user-study subset ($25$ of $30$ sessions). This suggests that the formula is not fragile around the chosen weights, but the claim is local: it does not cover arbitrary reweighting.

\paragraph{Agreement was checked on kinds, not ranking.}
The Cohen's $\kappa$ in \S\ref{sec:results-kappa} shows that three independent labellers usually agree on which divergence kinds appear in each session. It does not test whether they would order those divergence points in the same way as the score, even though the dashboard shows the highest-ranked points first. Testing that would need a separate ranking study, for example using Kendall $\tau$-b against the production ranking. This is left as future work in Chapter~\ref{ch:conclusion}.

\section{Study design and statistical power}
\label{sec:threats-study-design}

\paragraph{The user study is randomised, but too small for hypothesis testing.}
Five participants completed six sessions each, giving $30$ sessions in total. Group assignment was drawn uniformly at random before the study began, but the groups are still very small ($n = 3$ vs $n = 2$). Since the variation within each group is similar in size to the difference between groups, the chapter reports only the direction of the effect, with no $p$-value or confidence interval. The no-feedback group follows the same playbook, so its downward trend makes it unlikely that the tasks simply became easier over time. It does not remove two other explanations: participants who saw feedback may have learned what the dashboard rewards, or they may have become more familiar with the session format rather than better at refactoring.

\paragraph{The author is one of three raters.}
One of the three labellers was the author, so the agreement results could be affected by author judgement. To make this visible, \S\ref{sec:results-kappa} reports pairwise $\kappa$ rather than only a single aggregate score. The two other labellers were later assigned to the with-feedback group in the user study, but all labels were completed before they performed any study sessions. Their labels therefore could not have been influenced by their later behaviour in the experiment. A remaining concern is that two of the three with-feedback participants had already seen the divergence-point taxonomy through the labelling task. For this reason, \S\ref{sec:results-userstudy} reports per-participant trajectories so this can be inspected directly rather than hidden inside the group mean.

\section{Generalisation and implementation scope}
\label{sec:threats-external}

\paragraph{The injection sessions are controlled exercises.}
The $45$ injection sessions in the primary test set follow a scripted playbook, with one specific bad behaviour injected into each session. This makes them useful for checking whether the detector finds known divergence kinds, but they are not typical refactoring sessions. The precision and recall figures in \S\ref{sec:results-divergence} should therefore be read in that controlled setting, rather than as direct estimates for ordinary development work.

\paragraph{\texttt{SuboptimalOrdering} is based on the playbook author's judgement.}
The five ordering-injection sessions are labelled \texttt{SuboptimalOrdering} because the playbook author chose an order they considered worse than the available alternatives. The score formula only confirms this for $1$ of the $5$ sessions (\S\ref{sec:results-divergence}, session~$022$). The other $4$ have zero or negative magnitude under the production formula. The \emph{Ordering} precision and recall numbers should therefore be interpreted with this limitation in mind.

\paragraph{The agent extension is a feasibility experiment, not an agent evaluation.}
The eight-agent, $48$-session extension is used to see how detectors designed around human IDE sessions behave on a different kind of actor. It is not an evaluation of the tool as an agent-monitoring framework. Three of the four detector kinds fit agent traces poorly (\S\ref{sec:results-userstudy-agent-limits}). Agents have no IDE refactoring surface, so \emph{IDE-Replay} fires on every structural refactoring without offering a realistic fix. The six-checkpoint commit-gap threshold is rarely reached within a six-session arc, and the $2$-second edit-burst debouncer collapses each agent turn into one burst. For this reason, the agent result should not be read as showing that feedback had no effect on agent behaviour. It mainly shows that these detectors are not well matched to agent edit patterns.

\paragraph{The implementation is limited to Java sessions in IntelliJ.}
The plugin records activity through IntelliJ's PSI and VFS hooks, and the synthesiser is built around JDT. This means the evaluation does not cover other IDEs, other languages, or team workflows such as pull-request review. Supporting those settings would require additional engineering and validation, so Chapter~\ref{ch:conclusion} treats them as future work rather than simple configuration changes.

\paragraph{The reorder enumerator has a hard size budget.}
The \emph{Ordering} detector skips reorder windows with more than $7$ refactoring steps, because enumerating all valid orderings becomes expensive beyond that point (\S\ref{sec:methodology-reorder}). The recorded sessions rarely reach this limit. However, a session with a longer uninterrupted chain of refactorings could lose its largest windows from the \emph{Ordering} analysis.

% Backward-compatible anchors for cross-references from other chapters.
\label{sec:threats-internal}
\label{sec:threats-statistical}
\label{sec:threats-per-experiment}
\label{sec:threats-implementation}
\label{sec:threats-out-of-scope}
